\chapter{Results}
\label{chap:results}

\section{Build time against document size}

\sidenote{The raw timings live in \texttt{measurements.csv} under
\texttt{data/}. Edit that file and every figure in this chapter
follows.}%
\Cref{fig:latency} plots median build time against page count for the
three strategies. The data is read directly from
\texttt{data/measurements.csv}, so the figure updates whenever the
benchmark is re-run.

\begin{widefigure}
  \centering
  \begin{tikzpicture}
    \begin{axis}[
        width=0.94\fullwidth, height=62mm,
        xlabel={Document size (pages)},
        ylabel={Median build time (s)},
        xmin=15, xmax=165, ymin=0,
        grid=major,
        grid style={hairline,line width=0.4pt},
        axis line style={muted,line width=0.5pt},
        tick label style={font=\sffamily\scriptsize,color=muted},
        label style={font=\sffamily\small,color=ink},
        legend style={font=\sffamily\scriptsize,draw=hairline,
                      fill=white,at={(0.02,0.97)},anchor=north west,
                      cells={anchor=west}},
        every axis plot/.append style={line width=1.1pt,mark size=1.6pt},
      ]
      \addplot[accent,mark=*]        table[x=pages,y=baseline,   col sep=comma]{data/measurements.csv};
      \addplot[accentdark,mark=square*,dashed] table[x=pages,y=cached,col sep=comma]{data/measurements.csv};
      \addplot[muted,mark=triangle*] table[x=pages,y=incremental,col sep=comma]{data/measurements.csv};
      \legend{Baseline, Cached, Incremental}
    \end{axis}
  \end{tikzpicture}
  \caption[Median build time by strategy]{Median build time against
    document size. The baseline grows super-linearly once the document
    exceeds the worker's memory budget at roughly \num{120} pages;
    incremental builds stay almost flat.}
  \label{fig:latency}
\end{widefigure}

\section{Where the time goes}

Splitting a single \num{100}-page build into its phases
(\cref{fig:breakdown}) shows that the fixed cost $\alpha$ of
\cref{eq:model} is dominated by package initialisation, which is exactly
the part a warm cache removes.

\begin{widefigure}
  \centering
  \begin{subfigure}[t]{0.48\fullwidth}
    \centering
    \begin{tikzpicture}
      \begin{axis}[
          width=\textwidth, height=48mm,
          ybar, bar width=9pt,
          symbolic x coords={Init,Pass 1,Pass 2,Pass 3},
          xtick=data, ymin=0, ymax=4.2,
          ylabel={Time (s)},
          nodes near coords, nodes near coords style={font=\sffamily\tiny,color=muted},
          grid=major, grid style={hairline,line width=0.4pt},
          axis line style={muted,line width=0.5pt},
          tick label style={font=\sffamily\scriptsize,color=muted},
          label style={font=\sffamily\scriptsize,color=ink},
        ]
        \addplot[draw=accent,fill=accent!22] coordinates
          {(Init,3.61) (Pass 1,2.14) (Pass 2,1.18) (Pass 3,1.01)};
      \end{axis}
    \end{tikzpicture}
    \caption{Cold baseline.}
    \label{fig:breakdown-cold}
  \end{subfigure}
  \hfill
  \begin{subfigure}[t]{0.48\fullwidth}
    \centering
    \begin{tikzpicture}
      \begin{axis}[
          width=\textwidth, height=48mm,
          ybar, bar width=9pt,
          symbolic x coords={Init,Pass 1,Pass 2,Pass 3},
          xtick=data, ymin=0, ymax=4.2,
          ylabel={Time (s)},
          nodes near coords, nodes near coords style={font=\sffamily\tiny,color=muted},
          grid=major, grid style={hairline,line width=0.4pt},
          axis line style={muted,line width=0.5pt},
          tick label style={font=\sffamily\scriptsize,color=muted},
          label style={font=\sffamily\scriptsize,color=ink},
        ]
        \addplot[draw=accentdark,fill=accentdark!22] coordinates
          {(Init,0.34) (Pass 1,1.96) (Pass 2,0.68) (Pass 3,0.54)};
      \end{axis}
    \end{tikzpicture}
    \caption{Warm cache.}
    \label{fig:breakdown-warm}
  \end{subfigure}
  \caption[Phase breakdown]{Phase breakdown of a \num{100}-page build.
    The cache removes almost all of the initialisation cost but leaves
    the typesetting passes untouched.}
  \label{fig:breakdown}
\end{widefigure}

\section{Fitted parameters}

Fitting \cref{eq:model} to the measured series by ordinary least squares
gives the parameters in \cref{tab:fit}. The residual standard error is
below \qty{9}{\percent} for all three strategies, which is within the
run-to-run spread of the measurements themselves.

Two things stand out. First, the fixed cost $\alpha$ falls by more than
an order of magnitude as soon as the package tree is cached, confirming
the phase breakdown in \cref{fig:breakdown}. Second, the per-page cost
$\beta$ of the incremental strategy is not zero: even an unchanged page
must be re-shipped to the output driver, and that residual cost is what
keeps \cref{lem:bound} from being tight.

\begin{table}[htbp]
  \centering
  \caption{Least-squares fit of \cref{eq:model} to the measured series.}
  \label{tab:fit}
  \small
  \begin{tabular}{@{}l S[table-format=1.3] S[table-format=1.4] S[table-format=1.3] S[table-format=2.1]@{}}
    \toprule
    \textbf{\sffamily Strategy} & {$\alpha$ (s)} & {$\beta$ (s/page)} & {$\gamma$ (s)} & {RSE (\%)} \\
    \midrule
    Baseline    & 3.610 & 0.0742 & 0.284 & 8.4 \\
    Cached      & 0.340 & 0.0301 & 0.271 & 6.1 \\
    Incremental & 0.318 & 0.0061 & 0.119 & 4.7 \\
    \bottomrule
  \end{tabular}
\end{table}

\begin{keybox}[Crossover confirmed]
The measured crossover between the cached and incremental strategies
falls at $p^{\star} \approx 58$ pages, within \qty{4}{\percent} of the
value that \cref{thm:crossover} predicts from the fitted parameters in
\cref{tab:fit}.
\end{keybox}
